%% Formule de l'operateur : Synchro_spline
%% Dans ce fichier, les �ventuels accents sont cod�s
%% en ISO-8859-15 (aka latin9)
%% La ligne suivante donne trois exemples 
%% e accent aigu : �, e accent grave : �, e-dans-l'o : �

Les traitements suivants sont r�alis�s :
\begin{itemize}
\item  Dans un premier temps les valeurs pour lesquelles il y a un retour en arri�re dans la base de temps de $X$ sont supprim�es. \\
\item Puis, un d�calage temporel $D$ est appliqu� � $X$.\\
\item Enfin une interpolation par splines cubiques est faite avec la fonction Numpy  \textit{UnivariateSpline} (avec le param�tre s=0).\\
\end{itemize}

Si $Unite_X$ vaut "radians" ou "degres" alors ce sont $cos(X)$ et $sin(X)$ qui sont interpol�s. $X_{sync}$ est ensuite calcul� par l'arctangente (fonction \textit{arctan2 }de Numpy). \\

Cet op�rateur ne remplace pas  les valeurs manquantes dans $X$. C'est-�-dire, les valeurs qui sont manquantes dans $X$ sont aussi manquantes dans $X_{sync}$

